\section*{Capítulo \ref{ch:b3:nervous-systems} --- Organización de los sistemas nerviosos}

\begin{solution}{exo:b3:nervous-systems:1}
\omterm{def:b3:nervous-systems:cells}{Astrocitos}: amortiguan el potasio y el glutamato, alimentan a las
neuronas, inducen la \omterm{prop:b3:nervous-systems:barrier-budget}{barrera hematoencefálica} y regulan el flujo
sanguíneo local. \omterm{def:b3:nervous-systems:cells}{Oligodendrocitos}: mielinizan los axones del sistema
nervioso central. Células de Schwann: mielinizan los axones periféricos
(y guían su regeneración). \omterm{def:b3:nervous-systems:cells}{Microglía}: los \omterm{def:b3:innate-immunity:innate}{macrófagos} del encéfalo, que
podan sinapsis y retiran los restos. (Las células ependimarias tapizan
los ventrículos y fabrican el \omterm{prop:b3:nervous-systems:barrier-budget}{líquido cefalorraquídeo}.)
\end{solution}

\begin{solution}{exo:b3:nervous-systems:2}
La tinción ennegrecía por completo unas pocas neuronas y dejaba
invisibles las demás, de modo que podía verse entera una sola célula.
Cajal concluyó que las neuronas son células separadas que se contactan
sin fusionarse, con un sentido fijo de flujo de las dendritas al axón;
Golgi concluyó que formaban una red continua. Cajal tenía razón; el
microscopio electrónico mostró la hendidura sináptica en 1954.
\end{solution}

\begin{solution}{exo:b3:nervous-systems:3}
Un golpe estira el cuádriceps; sus husos musculares descargan; las
aferentes Ia entran en la médula y hacen sinapsis directamente sobre
las neuronas motoras del cuádriceps, que descargan y contraen el
músculo, mientras una \omterm{def:b3:nervous-systems:cells}{interneurona} inhibe las neuronas motoras de los
isquiotibiales. Una sinapsis y axones cortos a \qty{80}{m/s}: unos
\qty{30}{ms}. Una patada voluntaria pasa por el encéfalo ---percepción,
decisión, planificación motora, vías descendentes, varias sinapsis--- y
tarda \qtyrange{200}{300}{ms}.
\end{solution}

\begin{solution}{exo:b3:nervous-systems:4}
Simpático (noradrenalina en la diana): corazón más rápido y más fuerte,
pupila dilatada, motilidad y secreción intestinales reducidas, vías
respiratorias dilatadas. Parasimpático (acetilcolina): corazón
enlentecido, pupila contraída, motilidad y secreción intestinales
aumentadas, vías respiratorias contraídas.
\end{solution}

\begin{solution}{exo:b3:nervous-systems:5}
$\lambda = \sqrt{R_{m}d/4R_{i}} = \sqrt{10^{4}\times 10^{-4}/400}$ cm
$= \qty{0.05}{cm} = \qty{0.5}{mm}$. Tras \qty{2}{mm}: $e^{-4} \approx
2\,\%$. Duplicar $\lambda$ exige cuatro veces el diámetro,
\qty{4}{\micro m}.
\end{solution}

\begin{solution}{exo:b3:nervous-systems:6}
Sensitivo: $12\times 6 = \qty{72}{m/s}$, $1.2/72 = \qty{17}{ms}$. Motor:
\qty{90}{m/s}, $1.1/90 = \qty{12}{ms}$. Dos sinapsis, \qty{1}{ms}: unos
\qty{30}{ms}. Amielínicos: $1.1\sqrt{12} = \qty{3.8}{m/s}$ y
$1.1\sqrt{15} = \qty{4.3}{m/s}$: $315 + 258 + 1 \approx \qty{570}{ms}$
---demasiado lento para salvar el pie.
\end{solution}

\begin{solution}{exo:b3:nervous-systems:7}
$2.4\times 10^{20}/8.6\times 10^{10} \approx 2.8\times 10^{9}$ ATP por
neurona y segundo; a $2\times 10^{9}$ por espiga, unos \qty{1.4}{Hz} de
media. El código ha de ser escaso: pocas espigas, con la información
llevada por qué neuronas descargan y cuándo, no por frecuencias altas
en todas partes.
\end{solution}

\begin{solution}{exo:b3:nervous-systems:8}
Dos neuronas que se inhiben mutuamente sin cansarse forman un
conmutador: la primera que descarga suprime a la otra mientras dure el
impulso, y el relevo no llega nunca. La oscilación necesita un proceso
lento que debilite el dominio del vencedor ---adaptación de su
descarga, depresión de su sinapsis inhibidora o un rebote del perdedor
tras la liberación---, cuya constante de tiempo fija el período.
\end{solution}

\begin{solution}{exo:b3:nervous-systems:9}
La L-dopa atraviesa la barrera gracias al transportador de aminoácidos
neutros grandes; la dopamina, polar y cargada, no. Los pacientes toman
L-dopa (con un inhibidor de la enzima periférica para que no se
convierta antes de cruzar), y la propia enzima del encéfalo fabrica
dopamina a partir de ella allí donde hace falta. Los \omterm{def:b3:adaptive-immunity:antibody}{anticuerpos}, de
\qty{150}{kDa}, no cruzan ni las uniones estrechas ni transportador
alguno; su administración exige una lanzadera que secuestre un receptor
para la transcitosis, o la inyección en el \omterm{prop:b3:nervous-systems:barrier-budget}{líquido cefalorraquídeo}.
\end{solution}

\begin{solution}{exo:b3:nervous-systems:10}
La corriente de membrana por unidad de longitud es la resistiva
$V/r_{m}$ más la capacitiva $c_{m}\,\partial V/\partial t$; la
conservación de la carga da
$\partial I/\partial x = -(V/r_{m} + c_{m}\,\partial V/\partial t)$, y
con $\partial V/\partial x = -r_{i}I$, $\frac{1}{r_{i}}\,\partial^{2}V/
\partial x^{2} = V/r_{m} + c_{m}\,\partial V/\partial t$; al multiplicar
por $r_{m}$: $\lambda^{2}\,\partial^{2}V/\partial x^{2} = V + \tau\,\partial
V/\partial t$. Con $\partial V/\partial t = 0$ se recupera la ecuación
del teorema. $\lambda$ es una longitud y $\tau$ un tiempo, de modo que
$\lambda/\tau$ es una velocidad ---la escala natural de la rapidez con
que se propaga una perturbación.
\end{solution}

\begin{solution}{exo:b3:nervous-systems:11}
$v \propto \sqrt{d}$: cuadruplicar la velocidad exige dieciséis veces el
diámetro, \qty{8}{mm}. Área $\pi(4\,\text{mm})^{2} \approx \qty{50}{mm^2}$
frente a $\pi\,(10\,\unit{\micro m})^{2} = \qty{314}{\micro m^2}$: una razón de
$1.6\times 10^{5}$. Un nervio de fibras amielínicas rápidas es
imposible ---una sola de esas fibras sería tan gruesa como el nervio---,
de modo que un animal con muchos canales rápidos necesita \omterm{def:b3:nervous-systems:cells}{mielina}, y
por eso los vertebrados (y unos pocos invertebrados de forma
independiente) la desarrollaron.
\end{solution}

\begin{solution}{exo:b3:nervous-systems:12}
$a = 0.4/30^{0.75} = 0.4/12.8 = 0.031$; para \qty{70000}{g}: $0.031\times
\num{70000}^{0.75} = 0.031\times 4300 \approx \qty{134}{g}$. Los
\qty{1350}{g} humanos son diez veces la predicción. La masa del
encéfalo sigue a la masa corporal porque un cuerpo mayor necesita más
neuronas sensitivas y motoras, y porque las propias neuronas crecen con
el tamaño del encéfalo; lo que acompaña a la cognición es el número de
neuronas de la corteza, que depende de cuán densamente estén
empaquetadas ---los primates empaquetan más por gramo que los demás
mamíferos, y el ser humano más que ninguno.
\end{solution}

\begin{solution}{pb:b3:nervous-systems:1}
\textbf{1.} $\lambda = \sqrt{2\times 10^{4}\,d/400}$: para $d =
\qty{1e-4}{cm}$, $\sqrt{5\times 10^{-3}} = \qty{0.071}{cm} =
\qty{0.71}{mm}$; para \qty{4}{\micro m}, \qty{1.41}{mm}. $\tau = 2\times
10^{4}\times 10^{-6} = \qty{20}{ms}$.
\textbf{2.} $5\,e^{-0.3/0.71} = \qty{3.3}{mV}$; $5\,e^{-0.3/1.41} =
\qty{4.0}{mV}$.
\textbf{3.} Las entradas distales llegan atenuadas, de modo que el
lugar donde cae una sinapsis fija su peso; un potencial decae en
aproximadamente $\tau$, así que dos entradas solo se suman si caen a
unos \qty{20}{ms} una de otra ---la neurona es un detector de
coincidencias con una ventana de \qty{20}{ms}.
\textbf{4.} Mielínico, \qty{60}{m/s}; amielínico, $1.1\sqrt{10} =
\qty{3.5}{m/s}$; para alcanzar \qty{60}{m/s} sin \omterm{def:b3:nervous-systems:cells}{mielina}, $d = (60/1.1)^{2}
\approx \qty{3000}{\micro m}$: tres milímetros.
\textbf{5.} $\pi(5)^{2} = \qty{79}{\micro m^2}$ frente a $\pi(1500)^{2} =
7\times 10^{6}\,\unit{\micro m^2}$: unos \num{90000} axones mielínicos
en el espacio de uno solo.
\textbf{6.} \qty{2}{cm} a \qty{60}{m/s} son \qty{0.3}{ms}; a
\qty{3.5}{m/s}, \qty{5.7}{ms}: unos \qty{5}{ms} de retraso adicional.
Peor aún, la membrana desnuda tiene una capacitancia grande y pocos
canales, de modo que la corriente del último nódulo intacto puede no
llevarla al umbral: bloqueo de la conducción.
\textbf{7.} Sensitivo, $1.2/60 = \qty{20}{ms}$; dos sinapsis,
\qty{1}{ms}; motor, $1.1/90 = \qty{12}{ms}$: unos \qty{33}{ms} (más un
milisegundo en la unión neuromuscular).
\textbf{8.} Hasta el encéfalo: $0.6/60 = \qty{10}{ms}$ más dos
sinapsis, \qty{11}{ms} después de que la señal entre en la médula a los
\qty{20}{ms}: la señal llega a la corteza aproximadamente cuando el pie
se mueve; la experiencia consciente del dolor exige unos cientos de
milisegundos más de procesamiento, de modo que el pie se retira antes
de que se sienta el dolor.
\textbf{9.} $1.1\sqrt{1} = \qty{1.1}{m/s}$: $1.8/1.1 \approx
\qty{1.6}{s}$. El primer dolor, agudo, llega por fibras mielínicas
finas en decenas de milisegundos; el segundo dolor, sordo y urente, por
las fibras C amielínicas un segundo o más después.
\textbf{10.} Sensitivo, $3.2/60 = \qty{53}{ms}$; motor, $3.1/90 =
\qty{34}{ms}$; sinapsis, \qty{1}{ms}: unos \qty{90}{ms}. Los animales
largos tienen reflejos lentos; lo compensan con axones más gruesos y
delegando más en los circuitos locales.
\textbf{11.} $0.8/80 = \qty{10}{ms}$ en cada sentido, \qty{20}{ms};
sinapsis \qty{0.5}{ms}, unión \qty{1}{ms}, fuerza \qty{5}{ms}:
\qty{26.5}{ms}, y la transducción del propio huso y la conducción
dentro del músculo completan el resto de los \qty{30}{ms}.
\textbf{12.} Los axones amielínicos conducen a uno o dos metros por
segundo, de modo que los bucles que en un adulto tardan \qty{30}{ms}
tardan cientos, y con mala sincronía; la mielinización multiplica la
velocidad por diez o cincuenta y hace precisa la sincronía, que es lo
que exigen caminar, agarrar y hablar.
\textbf{13.} $20/50\,000 = 4\times 10^{-4}$ mol/s $= 2.4\times 10^{20}$
ATP por segundo; $2.8\times 10^{9}$ por neurona y segundo.
\textbf{14.} $4\times 10^{8} + 7000\times 2\times 10^{4} = 4\times 10^{8}
+ 1.4\times 10^{8} = 5.4\times 10^{8}$ ATP.
\textbf{15.} Todo en espigas: $2.8\times 10^{9}/5.4\times 10^{8} \approx
\qty{5}{Hz}$; la mitad: unos \qty{2.6}{Hz}.
\textbf{16.} A \qty{1}{Hz}, $5.4\times 10^{8}/2.8\times 10^{9} \approx
\qty{20}{\%}$ del presupuesto: el código es escaso ---la mayoría de las
neuronas callada la mayor parte del tiempo, y la información en cuáles
de las pocas que descargan lo hacen y cuándo.
\textbf{17.} $8.6\times 10^{10}\times 50\times 5.4\times 10^{8} = 2.3
\times 10^{21}$ ATP/s $= 3.9\times 10^{-3}$ mol/s $\approx \qty{190}{W}$
---el doble de todo el cuerpo en reposo.
\textbf{18.} $5\times 16 = \qty{80}{kJ/h} = \qty{22}{W}$: justo
suficiente, sin reserva alguna. Cuando el suministro se detiene, las
bombas se paran en segundos, los gradientes se agotan y la conciencia
se pierde en unos diez segundos.
\textbf{19.} $a = 0.4/12.8 = 0.031$; \qty{70}{kg}: $0.031\times 4300 =
\qty{134}{g}$; \qty{5000}{kg}: $0.031\times 1.06\times 10^{5} \approx
\qty{3300}{g}$.
\textbf{20.} Humano, $1350/134 \approx 10$; elefante, $4800/3300 \approx
1.5$.
\textbf{21.} En el \omterm{def:b3:nervous-systems:plan}{cerebelo} ---el \qty{98}{\%} de ellas---, coordinando
los cuarenta mil músculos de la trompa. Lo que acompaña a la cognición
es el número de neuronas corticales, no el total; el \omterm{def:b3:nervous-systems:plan}{cerebelo} acompaña
a la tarea motora de gobernar un cuerpo enorme.
\textbf{22.} Los axones y las dendritas se alargan con el encéfalo, de
modo que cada neurona ocupa más volumen y el número de neuronas crece
más despacio que la masa. En los primates el tamaño de las neuronas se
mantiene casi constante mientras el encéfalo crece, así que el número
de neuronas crece casi en proporción a la masa, y un encéfalo de
primate de un peso dado alberga varias veces las neuronas de un
encéfalo de roedor de ese mismo peso.
\textbf{23.} Un brazo puede coordinar sus propias ventosas y
articulaciones, alcanzar, agarrar, pasarse el alimento a lo largo de sí
mismo y evitar agarrar su propia piel, todo ello mediante circuitos
locales; no puede ver, ni elegir una diana, ni aprender una
discriminación. Prueba: sepárese el cordón nervioso del brazo del
encéfalo y tóquese el brazo con alimento ---lo agarra y lo desplaza
hacia donde estaría la boca; un brazo amputado hace lo mismo durante
una hora.
\textbf{24.} Unas $23$ sinapsis por neurona frente a $7000$:
trescientas veces menos. El mapa completo dice qué neurona puede
influir sobre cuál, pero no la fuerza ni el signo de cada conexión, ni
qué neuromoduladores las cambian, ni la dinámica ---de modo que incluso
en el gusano el diagrama de cableado predice la conducta solo en parte.
\textbf{25.} $\lambda = \qty{0.71}{mm}$ para la dendrita de
\qty{1}{\micro m}; latencia de retirada, unos \qty{33}{ms}; frecuencia
media de descarga, unos \qty{2.6}{Hz} con la mitad del presupuesto en
espigas.
\end{solution}
